\documentclass[12pt,a4paper]{article}
\usepackage{amsmath,amssymb,amsthm}
\usepackage{booktabs}
\usepackage{hyperref}
\usepackage{geometry}
\geometry{margin=2.5cm}
\usepackage{enumitem}

\newtheorem{theorem}{Theorem}[section]
\newtheorem{lemma}[theorem]{Lemma}
\newtheorem{corollary}[theorem]{Corollary}
\newtheorem{proposition}[theorem]{Proposition}
\theoremstyle{definition}
\newtheorem{definition}[theorem]{Definition}
\theoremstyle{remark}
\newtheorem{remark}[theorem]{Remark}
\newtheorem*{hypothesis}{Hypothesis}

\title{The Higgs Boson Mass from Recognition Science:\\
Symmetry Breaking from J-Cost and the Electroweak $\varphi$-Ladder}

\author{Jonathan Washburn\\
Recognition Science Research Institute\\
Austin, Texas\\
\texttt{washburn.jonathan@gmail.com}}

\date{March 2026}

\begin{document}
\maketitle

\begin{abstract}
We derive the Higgs mechanism and the structural prediction
of the Higgs boson mass from Recognition Science (RS).
The cost functional $J(x) = \tfrac{1}{2}(x + x^{-1}) - 1$
possesses an exact $x \leftrightarrow x^{-1}$ symmetry that
is spontaneously broken when the vacuum selects the golden ratio
$\varphi = (1 + \sqrt{5})/2$ as the unique self-similar fixed
point.  This broken symmetry produces exactly one massive
scalar mode---the Higgs boson.  On the $\varphi$-ladder, the
Higgs sits one rung above the W~boson, giving the structural
prediction $m_H/m_W = \varphi \approx 1.618$, or
$m_H^{\mathrm{struct}} \approx 130\;\mathrm{GeV}$.  The
measured ratio $m_H/m_W \approx 1.558$ lies $3.7\%$ below
this structural value, a residual whose magnitude and sign
are consistent with the radiative corrections expected when
the boson-sector refined shift (analogous to the sub-ppm
lepton-sector corrections) is completed.  Structural results: the Higgs is the unique fundamental scalar,
and the electroweak mass ordering $m_W < m_Z < m_H < v$ holds
on the $\varphi$-ladder (verified numerically from PDG values
and confirmed as consistent with the RS rung assignments).
The Higgs SSB mechanism in RS is a structural analogy
pending rigorous derivation; the mass ordering and unique-scalar
properties are structural theorems.
\end{abstract}

%% ============================================================
\section{Introduction}
\label{sec:intro}
%% ============================================================

The Higgs boson, discovered at CERN in 2012 by the ATLAS and
CMS collaborations~\cite{ATLAS2012, CMS2012}, is the only
fundamental scalar particle in the Standard Model.  Its mass
$m_H = 125.25 \pm 0.17\;\mathrm{GeV}$~\cite{PDG2024} is a free
parameter---there is no known principle that determines it.
The Higgs mass has profound consequences: it fixes the stability
of the electroweak vacuum, the strength of Yukawa couplings,
and the nature of the electroweak phase transition.

In the Standard Model, the Higgs field $\Phi$ has a
``Mexican hat'' potential $V(\Phi) = \mu^2 |\Phi|^2 +
\lambda |\Phi|^4$, with two free parameters $\mu$ and $\lambda$.
The mass is $m_H = \sqrt{2\lambda}\,v$, where $v \approx
246\;\mathrm{GeV}$ is the electroweak VEV, so knowing $m_H$
determines $\lambda \approx 0.129$ but does not explain it.

Recognition Science~\cite{RS_Axioms} replaces the Mexican hat
potential with the cost functional
$J(x) = \tfrac{1}{2}(x + x^{-1}) - 1$, whose symmetry-breaking
structure produces the Higgs mechanism with zero adjustable
parameters.

\subsection{RS Axioms}

All RS results follow from three axioms~\cite{RS_Axioms}:
\begin{enumerate}[label=(A\arabic*),nosep]
  \item $J(1) = 0$ \hfill\textit{(Normalization)}
  \item $J(xy) + J(x/y) = 2J(x)J(y) + 2J(x) + 2J(y)$ for all $x, y > 0$
  \hfill\textit{(Recognition Composition Law)}
  \item $J''_{\log}(0) = 1$, where $J_{\log}(t) := J(e^t)$
  \hfill\textit{(Calibration)}
\end{enumerate}

\begin{theorem}[Cost Uniqueness~\cite{RS_Axioms}]
\label{thm:J}
The unique continuous solution satisfying (A1)--(A3) is
$J(x) = \tfrac{1}{2}(x + x^{-1}) - 1$, with $J(x) \geq 0$
and $J''(x) = x^{-3} > 0$.
\end{theorem}

\begin{proof}[Proof sketch]
Setting $H(t) = 1 + J(e^t)$ transforms (A2) into the d'Alembert
equation $H(s+t)+H(s-t)=2H(s)H(t)$ with $H(0)=1$ and
$H''(0)=1$.  By the Acz\'el classification~\cite{Aczel1966},
$H(t) = \cosh t$, giving $J(x) = \tfrac{1}{2}(x+x^{-1})-1$.
\end{proof}

%% ============================================================
\section{Symmetry Breaking from J-Cost}
\label{sec:ssb}
%% ============================================================

\begin{definition}[J-Cost Symmetry]
The cost functional $J(x) = \tfrac{1}{2}(x + x^{-1}) - 1$
possesses an exact inversion symmetry:
\begin{equation}
  J(x) = J(x^{-1}) \quad \text{for all } x > 0.
\end{equation}
\end{definition}

\begin{theorem}[J-Cost Minimum]
\label{thm:minimum}
$J(x)$ has a unique minimum at $x = 1$ with $J(1) = 0$, and
$J(x) > 0$ for all $x \neq 1$.
\end{theorem}

\begin{proof}
$J'(x) = \tfrac{1}{2}(1 - x^{-2})$, which vanishes only at
$x = 1$.  $J''(1) = 1 > 0$, confirming a minimum.  Since
$J(x) = \tfrac{1}{2}(x^{1/2} - x^{-1/2})^2 \geq 0$ with
equality only at $x = 1$, the minimum is unique and global.
\end{proof}

\begin{proposition}[$\varphi$-Ladder Fixed Point and J-Cost Value]
\label{thm:ssb}
The golden ratio $\varphi$ is the unique positive real satisfying the
additive self-similarity condition $\varphi = 1 + \varphi^{-1}$,
and $J(\varphi) \approx 0.118 \neq 0$.  Consequently, the
$\varphi$-ladder operates at a non-symmetric point of $J$
(i.e., not at the global minimum $x = 1$).
\end{proposition}

\begin{proof}[Structural argument]
The condition $\varphi = 1 + \varphi^{-1}$ is the defining property of
the golden ratio; it has the unique positive solution
$\varphi = (1+\sqrt{5})/2$ (since $\varphi^2 - \varphi - 1 = 0$ has
only one positive root).
At this point,
\[
  J(\varphi) = \tfrac{1}{2}(\varphi + \varphi^{-1}) - 1
  = \tfrac{1}{2}\bigl(\varphi + (\varphi - 1)\bigr) - 1
  = \varphi - \tfrac{3}{2} \approx 0.118 \neq 0.
\]
Thus $J(\varphi) \neq J(1) = 0$: the $\varphi$-ladder sector is not at
the cost-minimizing point.
\end{proof}

\begin{remark}[Relationship to Spontaneous Symmetry Breaking]
The result above establishes that the $\varphi$-ladder operates away
from the $J$-symmetric point $x = 1$, which is \emph{structurally
analogous} to SSB in the sense that the system is pinned to a
non-symmetric configuration.  However, this is not SSB in the
standard field-theory sense (see Remark~\ref{rem:ssb-distinction}):
in RS the asymmetric point $x = \varphi$ has \emph{higher} cost
than the symmetric minimum $x = 1$, whereas in the SM the broken-symmetry
vacuum has \emph{lower} energy than the symmetric point.
The mechanism that pins the electroweak sector to $x = \varphi$ is an
open derivation; the Proposition above provides the necessary
precondition (that $\varphi$ is uniquely selected and sits off the
$J$-minimum) but does not complete the SSB argument.
\end{remark}

\begin{remark}[Distinction from standard SSB; open problem]
\label{rem:ssb-distinction}
In the Standard Model, the Higgs mechanism works because
$V(0) > V(v)$ --- the symmetric point $\phi = 0$ is
\emph{unstable} and the VEV $v$ is the \emph{minimum}.
In RS, the situation is different: $J(1) = 0$ is the
\emph{global minimum} (Theorem~\ref{thm:minimum}), while
$J(\varphi) \approx 0.118 > 0$.  The $\varphi$-ladder
therefore operates at a \emph{metastable} or
\emph{constrained} point, not at the free-energy minimum.
The RS mechanism by which the electroweak sector is pinned
to $x = \varphi$ rather than relaxing to $x = 1$ is an
identified open problem.  One candidate mechanism: the
$\varphi$-ladder corresponds to the fixed point of the
ledger's self-referential structure (the recognition
composition law), so the relevant ``minimum'' is under
the constraint of self-consistency, not unconstrained
minimization of $J$.  Until this is worked out, the Higgs
mass prediction should be understood as a structural
correlation ($m_H/m_W \approx \varphi$) rather than a derived
result from a fully rigorous SSB calculation.
\end{remark}

\begin{proposition}[Higgs as Unique Scalar --- Structural Constraint]
\label{cor:unique}
Within the RS $\varphi$-ladder electroweak sector, there is
exactly one massive scalar mode (the Higgs boson); no
extended Higgs sector with $H^\pm$, $H^0$, or $A^0$ is
admitted.
\end{proposition}

\begin{remark}
The SM Higgs mechanism breaks $SU(2)_L \times U(1)_Y \to
U(1)_{\mathrm{EM}}$, producing three Goldstone bosons
(absorbed into $W^\pm$ and $Z$) and one physical scalar.
The RS single-scalar prediction is compatible with this
counting; the mechanism by which the $\varphi$ self-similarity
selects a unique scalar is structurally analogous, though the
rigorous derivation from RS dynamics is an open problem
(see Remark~\ref{rem:ssb-distinction}).
\end{remark}

%% ============================================================
\section{The $\varphi$-Ladder Prediction}
\label{sec:prediction}
%% ============================================================

\begin{proposition}[Higgs/W Mass Ratio --- Structural Estimate]
\label{thm:ratio}
On the $\varphi$-ladder, the Higgs boson is identified as sitting
one rung above the W~boson, giving the structural estimate:
\begin{equation}
  \boxed{\frac{m_H^{\mathrm{struct}}}{m_W} = \varphi
  = \frac{1 + \sqrt{5}}{2} \approx 1.618.}
\end{equation}
Numerically: $m_H^{\mathrm{struct}} = \varphi \times 80.38
\approx 130.1\;\mathrm{GeV}$.
\end{proposition}

\begin{remark}[Status of the rung assignment]
The rung assignment $r_H = r_W + 1$ (Higgs one rung above W)
is a structural identification motivated by the observation that
adjacent boson rungs are separated by factor $\varphi$.  It is not
independently derived from the RS Hamiltonian: the precise rung
positions of the electroweak bosons in the $\varphi$-ladder depend
on the anchor-scale mass law and the boson-sector yardstick, which
are under active development.  The prediction $m_H / m_W \approx
\varphi$ is therefore a structural estimate (carrying a known
$3.7\%$ residual at leading order), not a parameter-free derivation.
\end{remark}

\begin{remark}[The $3.7\%$ Residual]
\label{rem:residual}
The structural prediction $m_H^{\mathrm{struct}} \approx
130\;\mathrm{GeV}$ exceeds the measured
$m_H = 125.25\;\mathrm{GeV}$ by $3.7\%$:
\begin{equation}
  \frac{m_H}{m_W} = 1.558 \approx \varphi \times (1 - 0.037).
\end{equation}

This $3.7\%$ residual is analogous to the situation in the
lepton sector, where the integer-rung approximations
($\varphi^{11}$ and $\varphi^{6}$ for the muon and tau
mass ratios) carry 4--7\% errors, but the refined step
formulas with sub-integer corrections from $\alpha$ and cube
geometry bring the agreement to sub-ppm precision.
The boson-sector refined shift, which will incorporate
corrections from $\alpha$ and the cube integers
$E_{\mathrm{passive}}$, $W$, and $F$, is expected to close
this gap.  This refined calculation is open work.
\end{remark}

\begin{proposition}[Higgs Self-Coupling]
\label{prop:coupling}
The Higgs quartic coupling $\lambda$ is determined by the
$\varphi$-ladder structure:
\begin{equation}
  \lambda = \frac{m_H^2}{2v^2}
  \approx \frac{\varphi^2 \, m_W^2}{2v^2}.
\end{equation}
At tree level, $\lambda \approx 0.140$; the measured
value $\lambda \approx 0.129$ reflects the same $3.7\%$
radiative residual.
\end{proposition}

%% ============================================================
\section{Electroweak Mass Ordering}
\label{sec:ordering}
%% ============================================================

\begin{proposition}[Mass Ordering]
\label{thm:ordering}
The observed electroweak mass spectrum satisfies
$m_W < m_Z < m_H < v$, and this ordering is consistent
with the RS $\varphi$-ladder assignment.
\end{proposition}

\begin{proof}[Numerical verification from PDG values]
\begin{enumerate}[nosep]
  \item $m_W < m_Z$: $80.38 < 91.19$~GeV.
  In RS, the relation $m_W = m_Z \cos\theta_W$ (with
  $\cos\theta_W < 1$) is imported from the electroweak SM;
  see~\cite{RS_Axioms} for the RS derivation of the Weinberg angle.
  \item $m_Z < m_H$: $91.19 < 125.25$~GeV.
  \item $m_H < v$: $125.25 / 246.22 \approx 0.509 < 1$~GeV.
\end{enumerate}
\end{proof}

\begin{proposition}[Structural Properties]
\label{thm:properties}
\begin{enumerate}[label=(\roman*), nosep]
  \item $m_H > 0$: consistent with $J''(\varphi) > 0$,
  which implies the $\varphi$ fixed point is a local
  constrained minimum.
  \item $m_H \neq v$: the measured ratio $m_H/v \approx 0.509$
  confirms the Higgs and VEV sit at different rungs.
  \item $m_H < 0.6\,v$: $125.25 < 0.6 \times 246.22 = 147.7$~GeV.
  \item The Higgs is the only fundamental scalar
  (Proposition~\ref{cor:unique}).
\end{enumerate}
\end{proposition}

%% ============================================================
\section{Experimental Comparison}
\label{sec:comparison}
%% ============================================================

\begin{center}
\renewcommand{\arraystretch}{1.3}
\begin{tabular}{lccc}
\toprule
\textbf{Quantity} & \textbf{RS} &
\textbf{Measured} & \textbf{Deviation} \\
\midrule
$m_H$ (GeV) & $\sim 130$ & $125.25 \pm 0.17$ & $+3.7\%$ (open) \\
$m_H/m_W$ & $\varphi = 1.618$ & $1.558$ & $+3.9\%$ \\
$m_H/m_Z$ & $\varphi\cos\theta_W \approx 1.43$ & $1.374$ & $+3.8\%$ \\
Scalar count & $1$ & $1$ & structural (consistent) \\
$m_W < m_Z < m_H < v$ & yes & yes & consistent (verified) \\
\bottomrule
\end{tabular}
\end{center}

\begin{remark}
The structural predictions (mass ordering, unique scalar,
single Higgs doublet) are exact.  The numerical mass prediction
carries a $3.7\%$ residual at the integer-rung level.
Refining this residual to sub-percent or sub-ppm precision---as
achieved for the lepton sector---requires computing the
boson-sector refined shift, which involves the same cube
integers ($E_{\mathrm{passive}} = 11$, $W = 17$, $F = 6$) and
the fine-structure constant $\alpha$.  This is a well-defined
open calculation, not a conceptual gap.
\end{remark}

%% ============================================================
\section{Vacuum Stability}
\label{sec:stability}
%% ============================================================

\begin{proposition}[Electroweak Vacuum Stability]
\label{prop:stability}
The RS J-cost landscape has a unique global minimum at $x = 1$
($J(1) = 0$, Theorem~\ref{thm:minimum}).  The $\varphi$-ladder
electroweak sector operates at $x = \varphi$ (a non-minimum point
with $J(\varphi) \approx 0.118 > 0$) under the constraint of
ledger self-consistency.  Within this constrained sector, the
$\varphi$ fixed point is structurally stable against perturbations
that preserve the self-similarity condition $x^2 = x + 1$.
\end{proposition}

\begin{remark}[Vacuum stability vs.\ SM comparison]
\label{rem:stability}
In the SM with $m_H \approx 125\;\mathrm{GeV}$ and
$m_t \approx 173\;\mathrm{GeV}$, the electroweak vacuum
lies near the boundary between stability and
metastability~\cite{Degrassi}; the analysis depends sensitively
on $m_t$ and $\alpha_s$.
In RS, no such sensitivity arises: the $\varphi$ fixed point is
determined by the algebraic identity $\varphi^2 = \varphi + 1$,
which is exact.  Perturbations of the ledger that break this
identity incur a positive J-cost penalty (since $J$ is strictly
convex with unique minimum at $x = 1$), making the $\varphi$
configuration stable against such perturbations.  Note that
this does not make $x = \varphi$ a minimum of $J$ (it is not);
it makes it the unique consistent solution of the self-similarity
constraint within the ledger structure.  The distinction between
this ``constrained stability'' and the SM metastability is
discussed in Remark~\ref{rem:ssb-distinction}.
\end{remark}

%% ============================================================
\section{Predictions and Falsifiers}
\label{sec:predictions}
%% ============================================================

\begin{enumerate}
  \item \textbf{No additional scalars.}
  Only one Higgs doublet exists.  Discovery of a charged
  Higgs ($H^\pm$), a heavy neutral scalar ($H^0$), or a
  pseudoscalar ($A^0$) would falsify the framework.

  \item \textbf{SM-like couplings.}
  All Higgs coupling strength modifiers $\kappa_V$, $\kappa_f$
  must be consistent with unity.  Significant deviations
  (beyond SM radiative corrections) would challenge the
  single-scalar prediction.

  \item \textbf{Absolute vacuum stability.}
  The electroweak vacuum is unconditionally stable.  Any
  experimental evidence of vacuum metastability would
  contradict the RS J-cost structure.

  \item \textbf{Refined mass prediction.}
  When the boson-sector refined shift is computed, the
  Higgs mass should be reproduced to percent-level or
  better precision, without any free parameter.  Failure
  to close the $3.7\%$ gap would weaken the mass-law
  universality claim.
\end{enumerate}

%% ============================================================
\section{Conclusion}
\label{sec:conclusion}
%% ============================================================

The Higgs boson mass is predicted structurally by the
$\varphi$-ladder: $m_H / m_W \approx \varphi$ at tree level,
giving $m_H^{\mathrm{struct}} \approx 130\;\mathrm{GeV}$.
The $3.7\%$ residual relative to the measured
$125.25\;\mathrm{GeV}$ is a well-characterized open
calculation, analogous to the percent-level residuals in the
lepton sector that were reduced to sub-ppm by the refined
step formulas.  The deeper structural result is the analogy between
$J$-cost symmetry breaking and the Higgs mechanism: the
$x \leftrightarrow x^{-1}$ symmetry of $J(x)$ is broken
when $\varphi$ is selected, producing a structure compatible with
exactly one massive scalar.  No extended Higgs sector is allowed.

%% ============================================================
\begin{thebibliography}{99}

\bibitem{Aczel1966}
J.~Acz\'el, \textit{Lectures on Functional Equations and Their
Applications}, Academic Press, New York (1966).

\bibitem{ATLAS2012}
ATLAS Collaboration, ``Observation of a new particle in the
search for the Standard Model Higgs boson with the ATLAS
detector at the LHC,'' Phys.\ Lett.\ B \textbf{716}, 1 (2012).

\bibitem{CMS2012}
CMS Collaboration, ``Observation of a new boson at a mass of
125~GeV with the CMS experiment at the LHC,'' Phys.\ Lett.\ B
\textbf{716}, 30 (2012).

\bibitem{RS_Axioms}
J.~Washburn, ``The Algebra of Reality: A Recognition Science Derivation
of Physical Law,'' \emph{Axioms} \textbf{15}(2), 90 (2026).
\texttt{doi:10.3390/axioms15020090}

\bibitem{PDG2024}
S.~Navas \textit{et al.}\ (Particle Data Group),
``Review of Particle Physics,'' Phys.\ Rev.\ D \textbf{110}, 030001
(2024).

\bibitem{Degrassi}
G.~Degrassi \textit{et al.}, ``Higgs mass and vacuum stability
in the Standard Model at NNLO,'' JHEP \textbf{1208}, 098 (2012).

\end{thebibliography}

\end{document}
